\documentclass[11pt,letterpaper]{article}

\usepackage[margin=0.82in,headheight=15pt,footskip=28pt]{geometry}
\usepackage{amsthm}
\usepackage{newtxtext,newtxmath}
\usepackage[T1]{fontenc}
\usepackage{microtype}
\usepackage{mathtools,bm}
\usepackage{booktabs,tabularx,array}
\usepackage{enumitem}
\usepackage{xcolor}
\usepackage[most]{tcolorbox}
\usepackage{titlesec}
\usepackage{fancyhdr}
\usepackage{hyperref}
\usepackage{lastpage}

\definecolor{onyx}{HTML}{17191C}
\definecolor{graphite}{HTML}{4B4F55}
\definecolor{rulegray}{HTML}{B8BBC0}
\definecolor{papergray}{HTML}{F3F3F1}
\definecolor{gold}{HTML}{8B6B2D}
\definecolor{deepblue}{HTML}{244B66}

\hypersetup{
  colorlinks=true,
  linkcolor=onyx,
  urlcolor=deepblue,
  pdftitle={Four-Phase Addressing and Eight-State Relations},
  pdfauthor={Anthony Vito Coppa},
  pdfsubject={Finite phase addressing, eight-state pairwise relations, graph factorization, and exact state evolution},
  pdfkeywords={phase addressing, occurrence, exterior algebra, K8, one-factorization, relative phase, exact computation}
}

\pagestyle{fancy}
\fancyhf{}
\fancyhead[L]{\small A. V. Coppa}
\fancyhead[R]{\small Four-Phase Addressing and Eight-State Relations}
\fancyfoot[L]{\small A. V. Coppa -- 5 August 2026}
\fancyfoot[R]{\small \thepage\ of \pageref{LastPage}}
\renewcommand{\headrulewidth}{0.35pt}
\renewcommand{\footrulewidth}{0pt}

\titleformat{\section}
  {\Large\bfseries\color{onyx}}
  {\thesection}{0.7em}{}
  [\vspace{0.2em}\color{rulegray}\titlerule]
\titleformat{\subsection}
  {\large\bfseries\color{onyx}}
  {\thesubsection}{0.65em}{}
\titlespacing*{\section}{0pt}{2.2ex plus .4ex minus .2ex}{1.2ex}
\titlespacing*{\subsection}{0pt}{1.8ex plus .3ex minus .2ex}{0.8ex}

\setlist[itemize]{leftmargin=1.4em,itemsep=0.22em,topsep=0.35em}
\setlist[enumerate]{leftmargin=1.55em,itemsep=0.22em,topsep=0.35em}
\setlength{\parindent}{0pt}
\setlength{\parskip}{0.52em}
\renewcommand{\arraystretch}{1.15}

\newtheoremstyle{relationstyle}
  {0.7em}{0.7em}{\itshape}{}%
  {\bfseries\color{onyx}}{.}{0.5em}{}
\theoremstyle{relationstyle}
\newtheorem{theorem}{Theorem}[section]
\newtheorem{proposition}[theorem]{Proposition}
\newtheorem{corollary}[theorem]{Corollary}
\newtheorem{definition}[theorem]{Definition}
\newtheorem{lemma}[theorem]{Lemma}
\newtheorem{remark}[theorem]{Remark}

\newcommand{\R}{\mathbb{R}}
\newcommand{\C}{\mathbb{C}}
\newcommand{\N}{\mathbb{N}}
\newcommand{\Z}{\mathbb{Z}}
\newcommand{\T}{\mathbb{T}}
\newcommand{\End}{\operatorname{End}}
\newcommand{\NF}{\operatorname{NF}}
\newcommand{\Eval}{\operatorname{Eval}}
\newcommand{\Phase}{\operatorname{Phase}_4}
\newcommand{\Cycle}{\operatorname{Cycle}_4}
\newcommand{\TimeRead}{\operatorname{TimeRead}_4}
\newcommand{\UnRead}{\operatorname{UnRead}_4}
\newcommand{\Scan}{\operatorname{Scan}_{28}}
\newcommand{\Pairs}{\operatorname{Pairs}}
\newcommand{\Ledger}{\mathsf{Ledger}}
\newcommand{\State}{\mathsf{State}}
\newcommand{\Step}{\mathsf{Step}}
\newcommand{\snoc}{\operatorname{snoc}}
\newcommand{\diagT}{\T_{\mathrm{diag}}}
\newcommand{\wedgeTwo}{\Lambda^2}
\newcommand{\Zfour}{\Z/4\Z}
\newcommand{\Zseven}{\Z/7\Z}
\newcommand{\Ztwentyeight}{\Z/28\Z}

\newtcolorbox{abstractbox}{
  colback=papergray,
  colframe=rulegray,
  boxrule=0.5pt,
  arc=1.2pt,
  left=8pt,right=8pt,top=7pt,bottom=7pt
}
\newtcolorbox{boundarybox}{
  colback=white,
  colframe=gold,
  boxrule=0.65pt,
  arc=1.2pt,
  left=8pt,right=8pt,top=7pt,bottom=7pt,
  title={\bfseries Scope and evidence boundary},
  coltitle=onyx,
  fonttitle=\normalsize
}
\newtcolorbox{resultbox}{
  colback=papergray,
  colframe=onyx,
  boxrule=0.7pt,
  arc=1.2pt,
  left=9pt,right=9pt,top=8pt,bottom=8pt
}

\begin{document}
\hypersetup{pageanchor=false}

\begin{titlepage}
\thispagestyle{empty}
\vspace*{0.09\textheight}

{\fontsize{28}{33}\selectfont\bfseries\color{onyx} Four-Phase Addressing and Eight-State Relations\par}
\vspace{0.75em}
{\fontsize{17}{22}\selectfont\bfseries Finite Pairwise Coordinates and Exact State Evolution\par}

\vspace{2.0em}
\color{gold}\rule{0.62\textwidth}{1.1pt}\color{black}
\vspace{1.9em}

\begin{abstractbox}
\textbf{Abstract.}
A cyclic phase can return while the occurrence carrying that phase remains new.  This paper gives an exact finite architecture for that distinction.  Euclidean division yields a bijection $\N\cong\{0,1,2,3\}\times\N$, separating four visible phase positions from an unbounded cycle address.  Two consecutive phase windows supply eight distinct addressed states.  Their complete unordered relation set has $\binom82=28$ elements, represented by the coefficients of an alternating two-form on $\R^8$.  Vanishing triangle residues are proved equivalent to the existence of vertex potentials, unique up to a common additive constant.  A seven-round one-factorization of $K_8$, combined with the Chinese remainder theorem $\Z/28\Z\cong\Z/7\Z\times\Z/4\Z$, gives an explicit bijective scanner of all twenty-eight relations without identifying scan order with the relation set.  For eight unit phases, quotienting by common phase gives $\T^8/\diagT\cong\T^7$; the twenty-eight pairwise phase reads are invariant under the diagonal action and satisfy exact multiplicative triangle compatibility.  Finally, these components are assembled into a state-and-transition contract in which operator state, phase, occurrence, channel state, pairwise relations, scan coordinate, and receipt remain separately typed.  The principal result is a finite relational machine whose phase may close, whose occurrence may continue, and whose complete pairwise structure remains exactly recoverable.
\end{abstractbox}

\vfill
{\large Anthony Vito Coppa\par}
{\large Independent Research\par}
{\large Draft 0.1 -- 5 August 2026\par}
\vspace{0.6em}
{\small Jurisdiction: finite combinatorics, exterior algebra, compact abelian groups, and exact state records.\par}
\end{titlepage}
\hypersetup{pageanchor=true}
\pagenumbering{arabic}

\tableofcontents
\clearpage

\section{Purpose, jurisdiction, and main statement}

Finite phase and event identity are often collapsed into one index.  A periodic display then appears to repeat an earlier state even when a new traversal has occurred.  The present construction keeps four objects distinct:
\begin{enumerate}
  \item the visible phase class;
  \item the unbounded occurrence address;
  \item the complete pairwise relation space on an eight-state window;
  \item a chosen order for reading those relations.
\end{enumerate}

The first three are intrinsic to the declared carriers.  The fourth is a scanner: useful and exact, but chosen.

\begin{boundarybox}
The following separations are part of the result.
\begin{itemize}
  \item Phase equality is not occurrence equality.
  \item Eight addressed states are not eight values of one untyped scalar.
  \item A pairwise relation is not the coordinate used to scan it.
  \item A diagonal zero of an alternating form is not an absent event or zero vector.
  \item Common phase is invisible to relative phase, but it is not absent from the underlying eight-channel state.
  \item Metric duration, physical time, and any synchronization between occurrence and scan order require additional data and are not supplied here.
\end{itemize}
\end{boundarybox}

Let $P_4=\{0,1,2,3\}$.  Let $E_8=\R^8$ with ordered basis $e_0,\ldots,e_7$, and let $\Delta_{28}=\wedgeTwo(E_8^*)$.

\begin{theorem}[A.0 -- Four-phase, eight-state relational architecture]\label{thm:main}
There exists an explicit tuple
\[
\mathsf{Rel}_{4,8}
=
\bigl(
P_4,\TimeRead,\UnRead,
E_8,\Delta_{28},
\{M_a\}_{a\in\Zseven},\Scan,
\T^8/\diagT,
\State,\Step
\bigr)
\]
carrying simultaneously:
\begin{enumerate}
  \item a bijection $\N\cong P_4\times\N$ with exact successor and carry;
  \item four-phase return with strictly advancing occurrence address;
  \item two consecutive phase windows giving eight distinct addressed states;
  \item all twenty-eight unordered and fifty-six oriented pairwise relations;
  \item additive triangle compatibility and potential reconstruction;
  \item a seven-round perfect-matching factorization of $K_8$ and a bijective $28$-coordinate scanner;
  \item seven independent relative phases and twenty-eight central-phase-invariant pairwise phase reads;
  \item an exact state contract in which operator state, phase, occurrence, channel payload, pairwise reads, scanner query, and receipt remain separately recoverable.
\end{enumerate}
\end{theorem}

The proof is assembled in Section~\ref{sec:proof-main}.  Sections~\ref{sec:phase}--\ref{sec:state} construct each component independently.

\section{Four-phase addressing and retained occurrence}\label{sec:phase}

\begin{definition}[A.1 -- Phase and cycle reads]
For $n\in\N$, define
\[
\Phase(n)=n\bmod4\in P_4,
\qquad
\Cycle(n)=\left\lfloor\frac n4\right\rfloor\in\N,
\]
and
\[
\TimeRead(n)=(\Phase(n),\Cycle(n)).
\]
For $(p,c)\in P_4\times\N$, define
\[
\UnRead(p,c)=4c+p.
\]
\end{definition}

\begin{theorem}[A.2 -- Euclidean address isomorphism]\label{thm:euclid}
The maps
\[
\TimeRead:\N\longrightarrow P_4\times\N,
\qquad
\UnRead:P_4\times\N\longrightarrow\N
\]
are mutually inverse bijections.
\end{theorem}

\begin{proof}
Euclidean division gives unique $c\in\N$ and $p\in P_4$ such that $n=4c+p$.  Therefore
\[
\UnRead(\TimeRead(n))=n
\]
and
\[
\TimeRead(\UnRead(p,c))=(p,c).
\]
\end{proof}

\begin{corollary}[A.3 -- Successor with carry]\label{cor:carry}
If $\TimeRead(n)=(p,c)$, then
\[
\TimeRead(n+1)=
\begin{cases}
(p+1,c),&0\le p<3,\\
(0,c+1),&p=3.
\end{cases}
\]
In particular,
\[
\Phase(n+4)=\Phase(n),
\qquad
\Cycle(n+4)=\Cycle(n)+1.
\]
\end{corollary}

Fix an opaque run identifier $u$ and define the occurrence address
\[
o_n=(u,n).
\]

\begin{theorem}[A.4 -- Non-erasing phase return]\label{thm:phase-return}
For every $n\in\N$,
\[
\Phase(n+4)=\Phase(n),
\qquad
o_{n+4}\ne o_n.
\]
If a receipt is appended at each successor step, then the receipt after four further steps has strictly greater length even though the visible phase has returned.
\end{theorem}

\begin{proof}
The phase identity is Corollary~\ref{cor:carry}.  Occurrence inequality follows from $n+4\ne n$.  Strict receipt growth is the append law for finite sequences.
\end{proof}

\section{The eight-state window}\label{sec:eight}

\begin{definition}[A.5 -- Two-window address]
Let $L_2=\{0,1\}$.  For $\ell\in L_2$ and $p\in P_4$, define
\[
i(\ell,p)=4\ell+p\in\{0,1,\ldots,7\}.
\]
\end{definition}

\begin{proposition}[A.6 -- Eight distinct addressed states]\label{prop:eight}
The map
\[
i:L_2\times P_4\longrightarrow\{0,1,\ldots,7\}
\]
is a bijection.  Thus two consecutive four-phase windows contain eight distinct addresses even though the phase representatives occur twice.
\end{proposition}

\begin{proof}
The values for $\ell=0$ are $0,1,2,3$ and those for $\ell=1$ are $4,5,6,7$.  These sets are disjoint and exhaust the codomain.
\end{proof}

The eight addresses are represented by the basis vectors $e_0,\ldots,e_7$ of $E_8$.  The basis is an address carrier.  It does not assert that the underlying states are physically orthogonal or that every application must use a Euclidean metric.

\section{Complete pairwise relations}\label{sec:relations}

\begin{definition}[A.7 -- Alternating relation record]
Let
\[
\Delta_{28}=\wedgeTwo(E_8^*).
\]
For coefficients $d_{ij}\in\R$ with $0\le i<j\le7$, define
\[
\Omega_d=\sum_{0\le i<j\le7}d_{ij}\,e^i\wedge e^j.
\]
Extend the coefficients by
\[
d_{ji}=-d_{ij},
\qquad
d_{ii}=0.
\]
\end{definition}

\begin{theorem}[A.8 -- Twenty-eight unordered relations]\label{thm:28}
The relation space has dimension
\[
\dim\Delta_{28}=\binom82=28.
\]
It supplies one independent signed coordinate for every unordered pair of addressed states, hence twenty-eight unordered and fifty-six oriented relations.  For all $x,y\in E_8$,
\[
\Omega_d(x,x)=0,
\qquad
\Omega_d(y,x)=-\Omega_d(x,y).
\]
\end{theorem}

\begin{proof}
A basis of $\Lambda^2(E_8^*)$ is
\[
\{e^i\wedge e^j:0\le i<j\le7\},
\]
whose cardinality is $\binom82=28$.  Alternation gives the diagonal zero and orientation reversal.  Each unordered pair admits two orientations.
\end{proof}

\begin{definition}[A.9 -- Triangle residue]
For distinct $i,j,k$, define the oriented triangle residue
\[
h_{ijk}=d_{ij}+d_{jk}+d_{ki}.
\]
\end{definition}

\begin{theorem}[A.10 -- Additive triangle-potential criterion]\label{thm:triangle}
The following are equivalent:
\begin{enumerate}
  \item $h_{ijk}=0$ for every oriented triangle $(i,j,k)$;
  \item there exist vertex values $s_0,\ldots,s_7\in\R$ such that
  \[
  d_{ij}=s_j-s_i
  \]
  for every ordered pair $(i,j)$.
\end{enumerate}
When they exist, the vertex values are unique up to addition of one common constant.
\end{theorem}

\begin{proof}
If $d_{ij}=s_j-s_i$, then the triangle sum telescopes.  Conversely, fix $s_0=0$ and define $s_i=d_{0i}$.  The vanishing residue on $(0,i,j)$ gives
\[
0=d_{0i}+d_{ij}+d_{j0}=s_i+d_{ij}-s_j,
\]
so $d_{ij}=s_j-s_i$.  If $s$ and $s'$ produce the same $d_{ij}$, then
\[
(s'_j-s_j)=(s'_i-s_i)
\]
for all $i,j$, hence the difference is constant.
\end{proof}

The residue $h_{ijk}$ is a finite compatibility test on a complete relation record.  It is not spacetime curvature and does not by itself define transport outside the eight-address carrier.

\section{A seven-round scanner of all relations}\label{sec:scanner}

Label the eight vertices by
\[
V_8=\{\infty\}\cup\Zseven.
\]
For each $a\in\Zseven$, define
\begin{equation}\label{eq:matching}
M_a=
\bigl\{\{\infty,a\}\bigr\}
\cup
\bigl\{\{a+j,a-j\}:j=1,2,3\bigr\},
\end{equation}
with finite labels reduced modulo $7$.

\begin{theorem}[A.11 -- One-factorization of $K_8$]\label{thm:factor}
Each $M_a$ is a perfect matching of $K_8$.  The seven matchings are pairwise edge-disjoint and partition the full edge set:
\[
E(K_8)=\bigsqcup_{a\in\Zseven}M_a.
\]
\end{theorem}

\begin{proof}
Within $M_a$, the finite endpoints are
\[
a,\ a\pm1,\ a\pm2,\ a\pm3,
\]
which are all seven residues modulo $7$, while $\infty$ occurs once.  Thus $M_a$ covers every vertex once.

Every edge $\{\infty,a\}$ occurs in exactly one matching.  For a finite edge $\{x,y\}$, the element
\[
a=\frac{x+y}{2}\in\Zseven
\]
is uniquely defined because $2$ is invertible modulo $7$.  The nonzero difference $x-y$ determines a unique $j\in\{1,2,3\}$ up to sign, so $\{x,y\}=\{a+j,a-j\}$ occurs exactly in $M_a$.  Hence the matchings are disjoint and contain all $7\cdot4=28$ edges.
\end{proof}

Define
\[
\kappa(a,0)=\{\infty,a\},
\qquad
\kappa(a,b)=\{a+b,a-b\}
\quad(b=1,2,3).
\]
Since $4$ and $7$ are coprime, the Chinese remainder theorem gives
\[
\chi:\Ztwentyeight\overset{\cong}{\longrightarrow}\Zseven\times\Zfour,
\qquad
\chi(q)=(q\bmod7,q\bmod4).
\]

\begin{corollary}[A.12 -- Bijective relation scanner]\label{cor:scanner}
The map
\[
\Scan=\kappa\circ\chi:\Ztwentyeight\overset{\cong}{\longrightarrow}E(K_8)
\]
is a bijection.
\end{corollary}

\begin{proof}
The map $\chi$ is a bijection by the Chinese remainder theorem, and $\kappa$ lists the four edges of each of the seven disjoint perfect matchings from Theorem~\ref{thm:factor}.
\end{proof}

\begin{remark}[A.13 -- Scan order is not occurrence order]
The scanner coordinate $q\in\Ztwentyeight$ orders the twenty-eight relations.  The occurrence coordinate $n\in\N$ orders traversals.  No equation $q=q(n)$ is supplied or implied.  A synchronization between them is an additional policy and requires its own witness.
\end{remark}

\section{Relative phase and multiplicative compatibility}\label{sec:phasequotient}

Let $\T=U(1)$.  The diagonal subgroup
\[
\diagT=\{(\lambda,\ldots,\lambda):\lambda\in\T\}
\]
acts on $\T^8$ by common phase multiplication.

\begin{theorem}[A.14 -- Seven independent relative phases]\label{thm:quotient}
There is a group isomorphism
\[
\T^8/\diagT\cong\T^7.
\]
One explicit isomorphism is
\[
[z_0,\ldots,z_7]
\longmapsto
(z_1\overline{z_0},\ldots,z_7\overline{z_0}).
\]
\end{theorem}

\begin{proof}
The displayed map is invariant under common multiplication by $\lambda\in\T$.  It is a homomorphism.  Every element of $\T^7$ is represented by $(1,w_1,\ldots,w_7)$, so the map is surjective.  Its kernel consists exactly of tuples with $z_j\overline{z_0}=1$ for all $j$, namely the diagonal subgroup.
\end{proof}

For $z=(z_0,\ldots,z_7)\in\T^8$, define the pairwise phase read
\[
w_{ij}(z)=z_j\overline{z_i}.
\]

\begin{theorem}[A.15 -- Pairwise phase compatibility]\label{thm:phasecompat}
The pairwise reads satisfy
\[
w_{ii}=1,
\qquad
w_{ji}=w_{ij}^{-1},
\qquad
w_{ij}w_{jk}=w_{ik},
\]
and therefore
\[
w_{ij}w_{jk}w_{ki}=1
\]
for every oriented triangle.  They are invariant under common phase:
\[
w_{ij}(\lambda z)=w_{ij}(z)
\quad(\lambda\in\T).
\]
Conversely, any collection $w_{ij}\in\T$ satisfying these identities arises from a phase tuple $z\in\T^8$, unique up to common phase.
\end{theorem}

\begin{proof}
The identities follow directly from $w_{ij}=z_j\overline{z_i}$.  For the converse, choose $z_0=1$ and $z_i=w_{0i}$.  Compatibility gives
\[
w_{ij}=w_{i0}w_{0j}=\overline{w_{0i}}w_{0j}=z_j\overline{z_i}.
\]
Any second realization differs from the first by multiplication of every $z_i$ by the same element of $\T$.
\end{proof}

The additive criterion of Theorem~\ref{thm:triangle} and the multiplicative criterion of Theorem~\ref{thm:phasecompat} are parallel finite compatibility statements.  They are not identified: one lives in an additive real carrier, the other in a compact multiplicative group.

\section{Exact state and transition contract}\label{sec:state}

This section assembles the finite architecture around two imported interfaces:
\begin{enumerate}
  \item an exact operator evaluator with a canonical normal form $\NF$;
  \item an append-only receipt operation whose prefix and final payload remain recoverable.
\end{enumerate}
The internal proofs of those interfaces are not repeated here; the present results depend only on their stated contracts.

Let $\mathcal A$ be an exact operator carrier with coordinate space $\mathcal C$ and mutually inverse maps
\[
\Eval:\mathcal C\longrightarrow\mathcal A,
\qquad
\NF:\mathcal A\longrightarrow\mathcal C.
\]
Let $\mathcal G\subset\mathcal A$ be a declared generator set.  Let $R_n\in\Ledger$ be a finite receipt archive.

\begin{definition}[A.16 -- Exact finite-relational state]\label{def:state}
At occurrence $n$, define
\[
X_n=(\alpha_n,p_n,c_n,z_n,R_n),
\]
where
\[
\alpha_n\in\mathcal C,
\qquad
(p_n,c_n)=\TimeRead(n),
\qquad
z_n\in\T^8.
\]
The complete pairwise phase record $W(z_n)=(w_{ij}(z_n))$ is derived from $z_n$.  A scanner query $q\in\Ztwentyeight$ is supplied separately and returns the selected unordered pair $\Scan(q)$.
\end{definition}

\begin{definition}[A.17 -- Typed transition]\label{def:transition}
Let $g_n\in\mathcal G$ be the declared operator move, and let
\[
F_{g_n}:\T^8\longrightarrow\T^8
\]
be a declared exact channel update.  Define
\[
\Step_{g_n}(X_n)=X_{n+1}
\]
by
\begin{align*}
\alpha_{n+1}&=\NF\bigl(\Eval(\alpha_n)g_n\bigr),\\
(p_{n+1},c_{n+1})&=\TimeRead(n+1),\\
z_{n+1}&=F_{g_n}(z_n),\\
R_{n+1}&=\snoc(R_n,r_{n+1}),
\end{align*}
where $r_{n+1}$ retains at least the move, unreduced expression, reduced operator coordinate, input and output addresses, channel update, evidence status, and occurrence address.
\end{definition}

\begin{proposition}[A.18 -- Recoverable transition offices]\label{prop:state}
Under Definitions~\ref{def:state} and \ref{def:transition}, each transition keeps separately recoverable:
\begin{enumerate}
  \item exact operator state $\alpha_n$;
  \item visible phase $p_n$;
  \item completed-cycle count $c_n$;
  \item eight-channel phase state $z_n$;
  \item twenty-eight derived pairwise reads $W(z_n)$;
  \item independently supplied scanner coordinate $q$;
  \item occurrence receipt $R_n$.
\end{enumerate}
No equality among their values authorizes substitution among their carriers.
\end{proposition}

\begin{proof}
Each field is assigned a distinct type and projection.  The pairwise reads are functions of $z_n$, the phase and cycle are functions of $n$, and the scanner is a function of $q$.  The receipt stores the transition occurrence rather than replacing any of these values.  Recoverability is therefore by the displayed projections and append contract.
\end{proof}

\begin{remark}[A.19 -- Equivariance requirement for relative-phase dynamics]
If pairwise relative phases are the authoritative relational output, a common-phase-insensitive update should satisfy
\[
F_g(\lambda z)=\lambda F_g(z)
\]
for all $\lambda\in\T$.  Then $F_g$ descends to $\T^8/\diagT$.  This equivariance is an adapter obligation, not a consequence of the finite address construction.
\end{remark}

\section{Proof of the main theorem}\label{sec:proof-main}

\begin{proof}[Proof of Theorem~\ref{thm:main}]
Item 1 is Theorem~\ref{thm:euclid} and Corollary~\ref{cor:carry}.  Item 2 is Theorem~\ref{thm:phase-return}.  Item 3 is Proposition~\ref{prop:eight}.  Item 4 is Theorem~\ref{thm:28}.  Item 5 is Theorem~\ref{thm:triangle}.  Item 6 is Theorem~\ref{thm:factor} and Corollary~\ref{cor:scanner}.  Item 7 is Theorems~\ref{thm:quotient} and \ref{thm:phasecompat}.  Item 8 is Definitions~\ref{def:state}--\ref{def:transition} and Proposition~\ref{prop:state}.  Every component is explicit, and the strict type separations are recorded in Sections~1 and 9.
\end{proof}

\section{Proved data and strict scope}\label{sec:scope}

The proved data may be summarized by
\[
\mathsf{FiniteRelationData}
=
\bigl(
P_4,\TimeRead,\UnRead,
E_8,\Delta_{28},
\{M_a\},\Scan,
\T^8/\diagT,
W,\State,\Step
\bigr).
\]
This is a register of distinct carriers and maps, not one opaque index.

The paper proves:
\begin{itemize}
  \item exact four-phase addressing with unbounded retained occurrence;
  \item eight distinct addresses across two phase windows;
  \item all twenty-eight unordered and fifty-six oriented relations;
  \item additive and multiplicative triangle compatibility;
  \item a complete seven-round factorization and bijective relation scanner;
  \item seven relative-phase degrees of freedom after quotient by common phase;
  \item an exact state-and-transition interface preserving operator, address, channel, relation, scan, and receipt offices.
\end{itemize}

It does not supply metric duration, a physical clock, a unique relation-scanning order, a synchronization between scan and occurrence, a physical interpretation of the eight channels, or a theorem forcing any additional cyclic register.  In particular, no $137$-address or $548$-address construction is used in the result.  Those may be introduced only as separately declared extensions.

\section{Computational verification}

The companion regression script checks the Euclidean address bijection and carry law over a declared finite range; uniqueness of the eight-state window; the $28/56$ pair count; exact alternating identities; additive triangle reconstruction; the complete $K_8$ one-factorization; the bijective $\Z/28\Z$ scanner; common-phase invariance and multiplicative triangle identities over exact finite root-of-unity samples; and a finite state-transition witness with strict receipt growth.  These computations are regression gates for the displayed formulas; the universal proofs remain the arguments above.

\begin{thebibliography}{9}
\small

\bibitem{BondyMurty}
J.~A. Bondy and U.~S.~R. Murty,
\emph{Graph Theory},
Springer, London, 2008.

\bibitem{Greub}
W.~H. Greub,
\emph{Multilinear Algebra},
2nd ed., Springer, New York, 1978.

\bibitem{CoppaAlgebra}
A.~V. Coppa,
\emph{A Complete $2\times2$ Operator Algebra for Recurrence, Phase, and Reversible Transport},
2026.

\bibitem{CoppaLift}
A.~V. Coppa,
\emph{Central Lift, Representation Difference, and Retained Occurrence},
2026.

\end{thebibliography}

\end{document}
